\documentclass{article}

\usepackage{amsmath,amssymb}
\usepackage{xurl}
\usepackage[colorlinks=true,linkcolor=blue,citecolor=blue,urlcolor=blue]{hyperref}

% Shared commands for notes in docs/paper-gaps/.

\DeclareUnicodeCharacter{2080}{\ifmmode{}_{0}\else\textsubscript{0}\fi}
\DeclareUnicodeCharacter{2081}{\ifmmode{}_{1}\else\textsubscript{1}\fi}
\DeclareUnicodeCharacter{2082}{\ifmmode{}_{2}\else\textsubscript{2}\fi}
\DeclareUnicodeCharacter{2099}{\ifmmode{}_{n}\else\textsubscript{n}\fi}

\newcommand{\Lean}{\textsc{Lean}}
\ExplSyntaxOn
\NewDocumentCommand{\leanid}{m}{
  \ifmmode
    \textsf{\detokenize{#1}}
  \else
    \group_begin:
    \urlstyle{sf}
    \tl_set:Nn \l_tmpa_tl {#1}
    \tl_replace_all:Nnn \l_tmpa_tl {\_} {_}
    \exp_args:NV \path \l_tmpa_tl
    \group_end:
  \fi
}
\ExplSyntaxOff
\providecommand{\tr}{\operatorname{tr}}
\newcommand{\tnleanIssueBase}{https://github.com/LionSR/TNLean/issues/}
\newcommand{\tnleanPrBase}{https://github.com/LionSR/TNLean/pull/}
\newcommand{\tnleanDocBase}{https://sirui-lu.com/TNLean/docs/}
\newcommand{\ghissue}[1]{\href{\tnleanIssueBase#1}{GitHub issue \##1}}
\newcommand{\ghpr}[1]{\href{\tnleanPrBase#1}{PR \##1}}
\newcommand{\doclink}[1]{\href{\tnleanDocBase#1.html}{\path{#1}}}

% Dirac notation and the identity operator, as in blueprint/src/macros/common.tex.
% \providecommand keeps note-local definitions authoritative.
\usepackage{dsfont}
\providecommand{\ket}[1]{|#1\rangle}
\providecommand{\bra}[1]{\langle #1|}
\providecommand{\braket}[2]{\langle #1 | #2 \rangle}
\providecommand{\ketbra}[2]{|#1\rangle\!\langle #2|}
\providecommand{\Id}{\mathds{1}}
\providecommand{\one}{\mathds{1}}
\providecommand{\C}{\mathbb{C}}
\providecommand{\MN}[1]{M_{#1}(\C)}
\providecommand{\MD}{M_D(\C)}

% Verdict marker: \gapnote{<kind>}{<status>}. Kind is the policy
% classification (clarification, local-correction, scope-restriction,
% unfaithful, false-source, open-gap); status is open, wip (elimination
% underway), resolved, or historical. Machine-read by the site index;
% prints a verdict line.
\newcommand{\gapnote}[2]{\par\noindent\textbf{Verdict:} #1 (#2).\par}


\title{Schuch--P\'erez-Garc{\'\i}a--Cirac SPT Classification: Continuous
    Tensor Realization of the Ground-State Path}
\author{}
\date{2026-11-04}

\begin{document}
\maketitle
\gapnote{scope-restriction}{open}

\section{What the source states}

Section~II.C of \cite{Schuch2011Phases} defines two Hamiltonians $H_0,H_1$ to
be in the same phase under a symmetry $G$ when, after a common blocking, there
are a phase gauge for $U^0_g$ and $U^1_g$, a representation
$U_g=U^0_g\oplus U^1_g\oplus U^{\mathrm{path}}_g$ on
$\mathcal H_0\oplus\mathcal H_1\oplus\mathcal H_{\mathrm{path}}$, and a
translation-invariant two-local path $h_\gamma$ with the prescribed endpoints,
$\|h_\gamma\|_{\mathrm{op}}\le1$, continuity in $\gamma$, a spectral gap above
the ground-state manifold bounded below independently of $N$ and $\gamma$, and
$[H_\gamma,U_g^{\otimes N}]=0$, the endpoint Hamiltonians being supported on
orthogonal summands.\footnote{Source traceability:
    \path{Papers/1010.3732/paper_v3.tex}, lines~404--466.}
For gapped Hamiltonians with exact MPS ground states it suffices to treat
parent Hamiltonians, all of which are interchangeable (lines~517--526).

Section~II.F treats injective MPS, that is, unique ground states, after the
standing reduction to standard form by blocking (lines~349--356 and
807--818). It states that the phases are exactly labelled by the cohomology
class of the virtual action $V_g$ in $U_g\mathcal P=\mathcal P(V_g\otimes\bar
V_g)$ (lines~832--867). Equality of phases is the direct-sum interpolation
(lines~870--921). Separation is proved only within MPS: ``there cannot be a
smooth MPS path connecting two such states'' (lines~923--933), with the
explicit caveat that the extension beyond MPS is only conjectured
(lines~949--952).

The separation hypothesis concerns the Hamiltonian path. The proof is to
``show that along any well-behaved path $H_\gamma$, $\mathcal P_\gamma$ needs
to change continuously'' (lines~938--941). It fixes $\gamma$, takes the unique
ground state $|\mu[\mathcal P_\gamma]\rangle$ of $H_\gamma$ with $\mathcal
P_\gamma$ in standard form (lines~957--960), proves continuity of the ground
state from the gap (lines~961--989), and concludes: ``since $|\mu[\mathcal
P_\gamma]\rangle$ is a polynomial in $\mathcal P_\gamma$, it follows that
$\mathcal P_\gamma$ can be chosen to be a continuous function of $\gamma$ as
well'' (lines~989--991). Appendix~C then works with continuous $\mathcal
P_\gamma$ in one virtual space whose support may grow (lines~2653--2717). The
summary at lines~1069--1072 speaks of ``a smooth gapped path $H_\gamma$''. The
symmetry-breaking subsection glosses a well-behaved path of MPS as one that
``yields a smooth and gapped path in the set of parent Hamiltonians''
(lines~1213--1216), and its separation argument again derives continuity of
the ground states from ``continuity and gappedness of $H_\gamma$''
(lines~1251--1264). Thus continuity of the tensors $\mathcal P_\gamma$ is a
conclusion of the source argument, not a hypothesis of the source theorem.

\section{The formal phase condition, clause by clause}

The relation \leanid{MPSTensor.IsSameExactMPSGappedPhase} consists of the
following clauses.
\begin{enumerate}
    \item A common blocking length $k>0$ for both endpoints. Source: ``there
          exists a finite $k$ such that after blocking $k$ sites''
          (lines~409--411) and ``this definition should be understood after an
          appropriate blocking of sites'' (lines~465--466). Matches.
    \item A unitary representation $U$ on a common space $\mathbb C^m$ and
          isometries $E_0,E_1$ of the blocked endpoint spaces with orthogonal
          ranges, intertwining $U$ with the blocked endpoint actions. The ranges
          and their orthogonal complement are invariant under the unitary $U$,
          so $U=U^0\oplus U^1\oplus U^{\mathrm{path}}$ as in lines~447--453,
          with orthogonality imposed even for equal representations
          (lines~460--465). Matches.
    \item Unit-modulus characters multiplying the blocked endpoint actions.
          Source: the phase gauge $U^p_g\leftrightarrow e^{i\phi^p_g}U^p_g$
          (lines~444--447), applied to the blocked system as lines~465--466
          require. Matches.
    \item Endpoint interactions equal to the canonical two-site parents of the
          blocked, included tensors. Source: parent Hamiltonians
          (lines~358--386) and their interchangeability (lines~517--526).
          Matches.
    \item A path of Hermitian two-site interactions with the prescribed
          endpoints, operator norm at most one, continuity, one gap $\delta>0$
          above the least spectral value for every parameter and every ring
          length $N\ge2$, and commutation of every ring Hamiltonian with
          $U_g^{\otimes N}$.\footnote{Formal declaration:
              \leanid{MPSTensor.SymmetricGappedInteractionPath}.}
          Source: items~1--4 of lines~415--428 and lines~447--453. Matches.
    \item At every parameter and ring length $N\ge2$, the least-energy
          eigenspace is the line of a nonzero periodic MPS vector, and the
          periodic rays have a one-site injective representative of positive
          bond dimension on the blocked chain. Source: separation within MPS
          (lines~923--933), unique ground state $|\mu[\mathcal P_\gamma]\rangle$
          with $\mathcal P_\gamma$ in standard form (lines~957--960), injective
          meaning unique ground state (lines~813--816). Matches; nonvanishing
          of the vector follows from the injective representative.
    \item A common ambient bond dimension $M$ bounding these representatives.
          Appendix~C and lines~989--991 place every $\mathcal P_\gamma$ in one
          virtual space; the source never lets the bond dimension of the path
          tensors grow without bound. We read this bound as part of the source
          setting. Given it, quantum Wielandt theory supplies one blocking
          after which every normal $\mathcal P_\gamma$ is one-site injective,
          so the common blocking of clause~6 is the standing standard-form
          convention of lines~349--356.
    \item A continuous family $\gamma\mapsto B(\gamma)$ of tensors of the
          ambient bond dimension whose periodic vectors span the least-energy
          lines.\footnote{Formal declaration:
              \leanid{MPSTensor.ExactMPSGroundPath}.}
          This is \emph{stronger} than the source. The source derives such a
          family from the gap (lines~938--941 and 961--991); it does not
          assume it.
\end{enumerate}

\section{The endpoint domain}

The classification theorems take one-site injective endpoints of positive
bond dimension.\footnote{Formal declarations:
    \leanid{MPSTensor.isSameExactMPSGappedPhase_iff_cohomologous} and
    \leanid{MPSTensor.isSameExactMPSGappedPhase_iff_cohomologous_of_character}.}
This is the source's standing convention, not a restriction. The source
assumes without loss of generality that every MPS is in standard form, ``this
might involve the blocking of sites of the original model''
(lines~349--356), and it classifies ``the symmetry of the blocked MPS''
(lines~816--818). The injective case of standard form is one-site
injectivity of the blocked tensor. A tensor that becomes injective only after
blocking $L$ sites enters through its $L$-blocked tensor and the blocked
physical action, exactly as in the source. Positive bond dimension excludes
only the zero tensor. The endpoint character is the phase gauge of
lines~444--447 and Appendix~B, lines~2596--2632; this was resolved in
\path{spc11_spt_interpolation_upper_range.tex}. The source labels phases by
the virtual action determined by lines~836--842; the theorems quantify over
any such prescribed action, whose existence for symmetric injective tensors
is proved separately.\footnote{Formal declaration:
    \leanid{MPSTensor.virtual_rep_of_symmetric_injective}.}
The common blocking of the phase condition is the source's ``after an
appropriate blocking''. No further statement for unblocked originals is
needed.

\section{Verdict}

The forward direction is faithful. The formal construction produces every
clause above, including the continuous tensor family, so it implies the
source relation.\footnote{Formal declaration:
    \leanid{MPSTensor.isSameExactMPSGappedPhase_of_isInjective_cohomologous}.}
The converse is restricted. It assumes, through clause~8, that the ground
lines of the gapped path have a continuous tensor realization in one bond
dimension. The source asserts that this realization exists along every
symmetric gapped path with unique MPS ground states. Its argument for that
step is the remark that the MPS vector is a polynomial in the tensor, which
gives continuity of the vector from the tensor but not the reverse. The
blueprint theorem \path{thm:spt_classification} therefore stays
\path{notready}. The explicit-domain theorems are recorded as
scope-restricted entries whose statements name the continuous exact MPS
phase condition.

\section{Elimination plan}

The source-faithful converse replaces clause~8 by the pointwise data of
clauses~6 and~7: for every $\gamma$, a one-site injective tensor $A_\gamma$ of
positive bond dimension at most $M$ whose periodic vectors span the
least-energy eigenspace of every ring Hamiltonian, with no continuity in
$\gamma$. The proof has three steps.
\begin{enumerate}
    \item Ground-line continuity. For a symmetric gapped interaction path
          whose least-energy eigenspace is one-dimensional at every parameter,
          show that its orthogonal projection is continuous in $\gamma$ for
          every fixed $N\ge2$. This is finite-dimensional: the uniform gap
          $\delta$ isolates the least eigenvalue, and the Riesz projection
          around it depends continuously on the continuous Hamiltonian. This
          is the source argument of lines~961--989.
    \item Replace ambient continuity by line convergence. The current converse
          uses continuity of the ambient family in two places: the common
          physical character\footnote{Formal declaration:
              \leanid{MPSTensor.ExactMPSGroundPath.exists_exact_symmetry}.}
          and the local class stability along convergent
          sequences.\footnote{Formal declaration:
              \leanid{MPSTensor.eventually_cohomologousTo_of_continuousAt_bounded_unital_periodic_rays}.}
          The character argument needs only a continuous invariant line at each
          ring length, which step~1 supplies. The class-stability argument
          needs convergence of normalized periodic vectors along a sequence;
          unit vectors in the ground lines, rephased against the limit, converge
          by step~1. The transfer bound from the physical gap is computed from
          the finite-ring ground lines and the pointwise minimal
          representatives, so it transfers unchanged once its inputs are
          restated.
    \item Restate the phase condition with the pointwise data and reprove the
          converse. The forward direction is unaffected, because its
          construction supplies the stronger continuous data.
\end{enumerate}
An alternative is to prove the source step literally: pointwise data with a
uniform bond-dimension bound admit a continuous ambient realization. That
statement needs a continuous choice through changes of minimal dimension,
which the source does not justify, and the plan above avoids it.

\bibliographystyle{alpha}
\bibliography{references}
\end{document}
